\section{The ten changes to the frozen files and rules}\label{app:changes}

This appendix carries, unchanged, the account of every change made to the frozen files and rules after the first freeze. It sits here rather than in Section 2.10 because a reader assessing the measurement needs the net effect, which Table 6 gives, and a reader auditing the measurement needs the detail, which is here. Nothing has been dropped in the move: the paragraphs below are the ones Section 2.10 carried in earlier versions of this paper, and a released check compares the two sets sentence by sentence.

Ten changes have been made to the frozen files and rules since, and each is recorded here with its measured effect rather than absorbed silently. The first is the Japanese extraction repair of Section 2.3, which was a defect in executing a stated exclusion rather than a change to the rule, and whose effect is given there. The second follows from it: the lexicon-validation script now drops a sample member that the repair removed from the authorising population, which changes the pooled false-negative rate from 11.1\% to 11.4\%, and the symmetrisation of Section 2.10 later moved it to 11.8\% by removing a second sample member from the population at risk. Weighting the three strata by their share of the tier-1 population, which the equal allocation requires and earlier versions omitted, moves the reported rate to 12.6\%. None of these touches the predicate.

The third is the correction of a defect in the determinacy lexicon that had been suppressing the clearest determinate construction in the English material, and it is the most consequential change in this paper. The pattern for the bare modal may carried a negative lookahead intended to withhold determinacy where a hedge follows, and the lookahead treated unless, without and only if alike. In a person may operate a standard remote identification unmanned aircraft only if the words only if are not a hedge on the permission but the determinate operator that conditions it, so the lookahead cancelled the very construction the lexicon exists to detect, and it did so wherever a permission was conditioned in the strongest form English offers. The repair removes only if, only when and only where from that lookahead and leaves unless and without in place. Its measured effect is to raise the returned count from 74 to 77 and the audited count from 26 to 27, leaving precision unchanged at that stage. The three provisions it adds are 14 CFR \S{}\S{} 107.110(b) and 107.115(b), adjudicated at this stage as false positives because the object of their condition is compliance with another section rather than an evaluation outcome, and 14 CFR \S{}89.110(b), which is a binding and which meets the defeater this design states below. The eighth change below reverses the first two: once a permission conditioned on meeting the requirements of another identified provision counts as reaching that provision's evaluation, both are bindings, and they are counted as such in all reported figures. We found the defect by taking that defeater seriously rather than by any check in our own verification suite, and that is itself a finding: the suite tests whether the reported numbers are mutually consistent, which they were throughout, and not whether the predicate detects the construction it claims to detect. A measurement can be internally consistent and wrong about its own object, and ours was.

The fourth is a repair to the penalty exclusion. Penalty provisions are barred from starting a path, but the pattern list carried only the Korean and Japanese forms, so three paragraphs of 14 CFR Part 89 imposing civil penalties for a false or concealed statement were admitted to the authorising denominator. The list now carries the United States forms as well. Its measured effect is to reduce the authorising count from 1,006 to 1,003, all three provisions falling in Part 89, and to leave the returned and audited counts untouched. The exclusion remains incomplete: the denominator audit of Section 2.8 drew three further penalty provisions in the Korean and Japanese material that the patterns do not catch, and we report that rather than repair it after the fact.

The fifth, sixth and seventh are extraction defects found after the GATE review of the previous version, and they are reported together because one mechanism produced two of them and because they changed the denominator. The fifth: the pattern that detects a top-level paragraph in the United States material matched any single lower-case letter in parentheses, so a romannumeral sub-item such as (i) was promoted to a paragraph of its own. That created three phantom authorising provisions, one of which, \S{}89.530(i), does not exist in 14 CFR and nonetheless appeared in the Part 89 tier distribution published here as the means by which a reader checks its own defeater. Worse, a promotion truncated its parent: \S{}\S{} 89.530(b) and (c) and \S{}\S{} 107.140(c)

and (d) each lost the text that followed the promoted item. The repair requires a top-level paragraph letter to be the successor of the previous one, which is how the Code of Federal Regulations is lettered.

The sixth: the Part 89 extractor had no fallback for a section consisting of a single unnumbered paragraph, so six sections entered the corpus as headings with no text. One of them was \S{}89.105, the operative prohibition of the whole Part, which provides that no person may operate an unmanned aircraft unless the operation meets the requirements of \S{}89.110 or \S{}89.115. A reviewer of the previous version observed that the paper never cites \S{}89.105; the reason was that it had never been measured. The repair gives that extractor the fallback its Part 107 counterpart already had.

The seventh: an article with no numbered paragraphs was held as a single node, so each EU definitions article counted as one authorising provision and fired on any term occurring anywhere in it, the largest being 17,101 characters and 68 definitions. The repair splits such an article into its numbered definitions, which also matches how EU instruments are cited. Two citations in Section 3.5 were wrong because of this defect and are corrected there.

Their combined measured effect is to take the corpus from 5,568 nodes and 8,596 edges to 5,758 and 8,834, the authorising count from 1,003 to 1,019, the returned count from 77 to 82 and the audited count from 27 to 30. Precision moves from 0.357 to 0.366. The surviving binding on the police route is unchanged.

The eighth is an adjudication rule rather than a file, and it is recorded here because it reversed two earlier labels. Where a permission is conditioned on meeting the requirements of another identified provision, and that provision names an evaluation, we now count the condition as reaching that evaluation. The alternative, requiring the evaluation to be named in the starting provision's own text, contradicts the path-based construct of Section 2.1 and would make the hop bound pointless; it also split two provisions of identical form, counting Article 5(1) of Implementing Regulation (EU) 2019/947 while rejecting 14 CFR \S{}107.110(b). A sweep of all 57 rejected provisions found the form in exactly three, so the change does not cascade. Against it we adopted a second rule, in the opposite direction: where the condition can be satisfied by more than one branch, it fails if any reachable branch leaves the permission standing without an evaluation. That rule costs us \S{}89.105(a), which would otherwise have been a second surviving binding on the police route. Both rules are applied to every returned provision, not only to the new ones, and the audit record names the provisions each one moved.

The ninth repair is of a different kind and was found by sampling rather than by review. The recall claim of Section 2.6 rests on the tier-1 population, and tier 0 had never been sampled at all, so we drew from it: a census of the 55 tier-0 authorising provisions carrying a determinate operator, and a random sample of 40 of the remaining 402, seed 20261009. The census found no missed binding of the kind the sweep looks for. It found something else. The lexicons are not symmetric across the two East Asian languages of the corpus. eval carried the Korean word for examination but not its Japanese counterpart, and det carried the Korean form for shall revoke but not the Japanese one, with three further pairs in the same state. We therefore added the Japanese counterpart of every concept the Korean lexicons already carried, in both directions: the hedge lexicon was symmetrised with the others, so the repair removes returns as well as adding them. This is a correction to coverage, not an extension of the construct; no concept is counted now that was not already counted in Korean.

Its measured effect is the largest of the ten. The returned count moves from 82 to 91, ten provisions entering and one leaving, and the audit of those ten finds eight bindings and two false positives. All eight are Japanese driving-licence provisions that condition the grant, the eligibility or the revocation of a licence on a named examination: Article 89(1), which requires the applicant to sit the driving licence examination; Article 104-2-2(1), (2) and (4), which require the public safety commission to revoke a licence where the holder fails or does not sit the re-examination; Article 90(13) and Article 90-2(2), which turn on having passed the examination; and Article 96(4) and (5), which condition eligibility to sit one examination on holding a licence obtained through another. The earlier version of this paragraph reported four bindings and six rejections here, and the blind re-adjudication of the audit moves four of those six the other way. The audited count moves from 33 to 41, precision from 0.402 to 0.451, and the ground rate from 7 of 604 to 15 of 604. The aviation-to-ground ratio falls from 9.12 to 4.26 and the Fisher test from p = 2.0 $\times$ 10$-$9 to p = 3.2 $\times$ 10$-$6 , so the contrast survives at a materially smaller effect size. The surviving binding on the police route, \S{}89.110(b), is unchanged, as is the carve-out result. We report the repair rather than the omission because its direction is against us, and a rule applied only when it helps is not a rule.

The tenth repair is the one a reviewer found, and it is the only one in this list that the study had already declared and then failed to apply. Specification F-5 requires an authorising provision to be a paragraph node, on the ground that an article caption carries no normative text. The implementation tested whether the provision identifier contained a full stop. For the Korean and Japanese instruments that is the right test, because an article is 43 and a paragraph is 43.1. For the United States instruments it is not, because a section identifier already contains a full stop: 107.13 is the section and 107.13(a) the paragraph. The rule was therefore never enforced on Part 107 or Part 89, and five section captions with paragraph children of their own stood in the authorising population: 14 CFR \S{}\S{} 107.13, 107.69, 89.215, 89.510 and 89.515. The repair replaces the full-stop test with the condition the specification actually states, which is that the node has no child node in the graph; that test does not depend on an identifier convention and applies to all sixteen instruments alike.

Its measured effect is to take the authorising count from 1,019 to 1,014 and the tier 1 and tier 0 counts from 494 and 434 to 491 and 432. The returned count, the audited count and the precision are unchanged at 91, 41 and 0.451, because no removed caption was returned. The aviation denominator falls from 246 to 241, so the aviation rate rises from 10.6\% to 10.8\%, the aviation-toground ratio from 4.26 to 4.34, Cohen's h from 0.346 to 0.353 and the aviation-against-ground Fisher test from p = 3.2 $\times$ 10$-$6 to p = 2.6 $\times$ 10$-$6 . The surviving binding on the police route and the carve-out result are unchanged. The direction of this repair is in our favour, which is why we state what it cost us to have missed it: the Part 89 distribution published in Section 3.4 as a check on the defeater was wrong in the only part a reader could have used, naming four tier-1 provisions when two of the four were captions.

\section{What the verification apparatus covers, and what it does not}\label{app:gates}

Eleven gates run over the consistency suite. One of them measures how weak another is, and one checks that the released archive contains what the data availability statement says it contains. The suite and the gates can establish that the numbers in this paper agree with the released outputs and with each other; they cannot establish that the outputs are correct, and the distinction matters enough to state the coverage of each gate rather than its pass count. The repository file \texttt{REPRODUCIBILITY.md} lists every gate with the inputs that break it, the inputs that do not, and the measured size of the gap where one has been measured. Four statements about coverage belong here, because earlier versions of this paper implied more than the apparatus delivers and because the apparatus itself has since been rebuilt.

The first concerns provenance. The gate that requires every measured quantity in the manuscript to be present in a released output compares values, not provenance. Its vocabulary is built by rounding every number in every released output to one, two and three decimal places and converting proportions to percentages, so a fabricated value that happens to coincide with any released number at the precision the manuscript prints it will pass. We measured the size of this by generating every decimal between zero and one at each precision and counting how many the vocabulary already holds: 10 of the 10 one-decimal values, 54 of the 100 two-decimal values, and 75 of the 1,000 three-decimal values. At one decimal place the vocabulary is saturated and the gate cannot fail. At two places a number has to clear a bar that about half the possible values do not clear, and at three places the bar is real. These three counts are a re-measurement. The script that produced the figures an earlier version printed was lost with the working tree, and the reimplementation written from the sentence above does not reproduce them; the released counts are the ones this paper now stands on, they are written by \texttt{verify/vocab}\\texttt{\_saturation}.json, and a released check fails if this paragraph and that output disagree. The direction of the correction is against us at one decimal place, where the earlier figure implied a bar that does not exist. A zero count from the provenance gate means no manuscript number failed to coincide with a released one; it does not mean every manuscript number came from the output a reader would expect it to come from.

The second concerns what replaced the apparatus this appendix used to describe. The consistency suite of the earlier version, and three gates beside it, were lost with the working tree and are not recovered. What runs now was built to close the defect classes this project actually produced rather than to restore a pass count, and its first design rule is that no check may take its expected value from the thing it checks. The gate that was previously reported as true but uninformative is the reason for that rule: the runner recompiled the typeset files immediately before checking that they were current, so inside the run the check could not fail. Expected values now come from the measurement or from the freeze record and never from what the build writes.

The third concerns negative checks. A check that a superseded number is absent passes as soon as that number is gone, which is not the same as the manuscript being right. Twice in this project a count was corrected while a statistic computed from it was left alone, and the absence check passed over it. The suite therefore pairs every negative check with a positive one that requires the current measured value to be present in the manuscript, and the figures that carry the counts are regenerated from the measurement rather than written by hand.

The fourth concerns what no check in this apparatus reads. Three defects in this manuscript were visible only in the typeset output and not in any source file: every table was printed twice, the text layer of four superseded figures was set as body prose, and the four figures were replaced by empty placeholder boxes by the assembler on every build. The suite now reads the typeset PDF directly for all three. That it had to be taught to do so is the general point of this appendix: an apparatus tests what it was pointed at, and a pass count says nothing about where it was pointed.
